\section{Threat Model and Goals}
\label{sec:threat}

\paragraph{System.} An operator runs a \mc-style cluster: disaggregated
prefill and decode pools, a pooled DRAM/SSD KV store with content-keyed
deduplication, and a cache-aware global scheduler. Clients are
\emph{principals}---API keys or tenants---that the operator authenticates. The
operator, the serving software and the hardware are trusted.

\paragraph{Attacker.} The attacker controls one or more principals, possibly
many (Sybils). It knows the model, block size, scheduling algorithm and
approximate hardware, and may submit arbitrary prompts at chosen times subject
to rate limits. For each request it observes everything a client sees: the
admission outcome (including HTTP 429), the arrival time of every output token
measured at its client, and the usage record (prompt, cached and output token
counts, and the resulting charge). Its goals are to learn whether a prefix was
processed for another principal (\emph{membership}), to reconstruct another
principal's prompt (\emph{reconstruction}), or to learn how recently or how
widely a prefix is used (\emph{activity}).

\paragraph{Out of scope.} Attackers with code on cluster hosts or GPUs
(microarchitectural channels, Rowhammer on KV memory~\cite{bitflip2026kv}),
network observers of encrypted token streams, a malicious operator,
semantic caches, non-prefix chunk reuse~\cite{yao2025cacheblend}, information
the model itself memorizes, and denial of service. Contention in the
\emph{decode} batch affects TBT in any shared system, isolated caches included;
we treat it as part of the public load signal and do not claim to remove it.

\subsection{Observations and the public load signal}
For a request $r$ from principal $p$, let
$o(r) = (a_r, \tau_r, \mathbf{t}_r, u_r)$ be its admission outcome, first-token
release time relative to arrival, later token times, and usage record. The
attacker's view is the sequence $\Obs_p$ of $o(r)$ over its requests.

Some dependence on others is unavoidable: if the cluster is busy, everyone
waits longer. \sys makes this dependence explicit through a \emph{public load
signal} $\Lpub$: a load level drawn from $k$ values, published once per epoch of
length $\Delta$. Every scheduling quantity that a client could observe may depend
on $\Lpub$, and on nothing else outside the principal's own requests.

\begin{definition}[Observational isolation]
\label{def:oi}
A serving system is \emph{observationally isolated} relative to $\Lpub$ if, for
every principal $p$ and any two executions in which $p$ submits the same
requests at the same times and $\Lpub$ takes the same values, the deterministic
part of $\Obs_p$ is identical, and the remaining randomness in $\Obs_p$ is
independent of every other principal's requests.
\end{definition}

An isolated cache with a fixed per-principal quota meets this definition but
loses all cross-principal reuse. Through $\Lpub$, the definition admits leakage
of at most $\log_2 k$ bits per epoch about \emph{aggregate} load, never about a
specific prefix.

\subsection{Goals}
\begin{itemize}
\item[\textbf{G1}] \textbf{Confidentiality.} Observational isolation for
  admission, first-token time and usage, with every violation (a missed release
  deadline, \cref{sec:design}) detected, counted and bounded.
\item[\textbf{G2}] \textbf{Preserved sharing.} Physical reuse equal to
  unrestricted global sharing: no block is recomputed because of \sys.
\item[\textbf{G3}] \textbf{Capacity.} More SLO-compliant requests than
  per-principal isolation on the same hardware.
\item[\textbf{G4}] \textbf{Deployability.} No change to the model, kernels or
  KV store; changes confined to the scheduler and the API front end.
\end{itemize}
